\documentclass[11pt,a4paper]{article}
\usepackage[utf8]{inputenc}
\usepackage[T1]{fontenc}
\usepackage[french]{babel}
\usepackage{lmodern}
\usepackage{geometry}
\usepackage{setspace}
\usepackage{booktabs}
\usepackage{hyperref}
\geometry{margin=2.4cm}
\onehalfspacing

\title{Rapport Scientifique Final --- Suite E00$\rightarrow$E11}
\author{Equipe Projet DAAA}
\date{\today}

\begin{document}
\maketitle
\tableofcontents
\newpage

\section{Introduction}
Ce rapport presente une etude experimentale progressive sur un encodeur Transformer audio pre-entraine par MAE puis adapte en ASR via CTC. 
L'objectif scientifique est triple: (i) verifier l'apport du pre-entrainement MAE, (ii) identifier les compromis entre WER, temps d'inference et memoire GPU, (iii) produire un benchmark multi-variantes avec consolidation statistique finale.

Contraintes de conception: environnement vdigpu, budget stockage limite (environ 15 Go), execution reproductible, et suivi frugal des ressources.

\paragraph{Choix globaux assumes}
\begin{itemize}
\item Screening uniquement sur LibriSpeech pour tenir le budget disque/temps.
\item VoxPopuli est exclu par defaut de cette suite (extension possible si marge).
\item Le cache data est conserve entre experiences.
\end{itemize}


\section{Methodologie generale}
\subsection{Pipeline}
Le protocole suit les etapes \texttt{make data}, \texttt{make train}, \texttt{make test}. 
Chaque experience est executee dans un namespace dedie, avec nettoyage des checkpoints intermediaires apres archivage des resultats, tout en conservant le cache de donnees.

\subsection{Metriques}
Les metriques principales sont:
\begin{itemize}
\item WER (Word Error Rate) pour la performance ASR.
\item Temps d'inference et debit (samples/s) pour la frugalite temporelle.
\item Memoire GPU pic pour la frugalite materielle.
\item Taille modele (\#parametres) pour l'analyse capacite/cout.
\end{itemize}

\subsection{Regles statistiques}
Le screening initial est realise en 1 seed par variante (E01--E08). 
La consolidation finale est realisee sur 5 seeds pour les 3 meilleures variantes (E09--E11), avec reporting moyenne~$\pm$~ecart-type.


\section{Experience E00 --- Smoke test pipeline}
\subsection{Choix et justification}
\textbf{Phase}: smoke\\
\textbf{Seeds prevues}: 42

Valider la chaine complete data/train/test avec un cout minimal avant toute analyse scientifique.

\paragraph{Hypothese attendue}
Aucune conclusion de performance; uniquement validation technique et reproductibilite.

\subsection{Resultats}
\textit{A completer apres execution des runs.}

\subsection{Discussion des resultats}
\textit{A completer apres analyse des resultats de cette experience.}


\section{Experience E01 --- Baseline small MAE}
\subsection{Choix et justification}
\textbf{Phase}: screening\\
\textbf{Seeds prevues}: 42

Etablir un point de reference frugal, stable et compatible avec la contrainte 15 Go.

\paragraph{Hypothese attendue}
Servir de base comparative pour toutes les autres variantes.

\subsection{Resultats}
\textit{A completer apres execution des runs.}

\subsection{Discussion des resultats}
\textit{A completer apres analyse des resultats de cette experience.}


\section{Experience E02 --- Small NoMAE}
\subsection{Choix et justification}
\textbf{Phase}: screening\\
\textbf{Seeds prevues}: 42

Isoler l'effet du pre-entrainement MAE en supprimant cette phase.

\paragraph{Hypothese attendue}
WER degrade vs E01 si MAE apporte une representation utile.

\subsection{Resultats}
\textit{A completer apres execution des runs.}

\subsection{Discussion des resultats}
\textit{A completer apres analyse des resultats de cette experience.}


\section{Experience E03 --- Tiny MAE}
\subsection{Choix et justification}
\textbf{Phase}: screening\\
\textbf{Seeds prevues}: 42

Tester une architecture plus compacte pour maximiser la frugalite.

\paragraph{Hypothese attendue}
Moins de memoire/temps, possible degradation WER.

\subsection{Resultats}
\textit{A completer apres execution des runs.}

\subsection{Discussion des resultats}
\textit{A completer apres analyse des resultats de cette experience.}


\section{Experience E04 --- Base MAE}
\subsection{Choix et justification}
\textbf{Phase}: screening\\
\textbf{Seeds prevues}: 42

Tester une capacite superieure pour evaluer le compromis performance/cout.

\paragraph{Hypothese attendue}
WER potentiellement meilleur, mais cout memoire/temps plus eleve.

\subsection{Resultats}
\textit{A completer apres execution des runs.}

\subsection{Discussion des resultats}
\textit{A completer apres analyse des resultats de cette experience.}


\section{Experience E05 --- Small Learned Positional Embeddings}
\subsection{Choix et justification}
\textbf{Phase}: screening\\
\textbf{Seeds prevues}: 42

Comparer sinusoidal vs learned embeddings en gardant le reste fixe.

\paragraph{Hypothese attendue}
Impact sur la qualite de sequence modeling et le transfert ASR.

\subsection{Resultats}
\textit{A completer apres execution des runs.}

\subsection{Discussion des resultats}
\textit{A completer apres analyse des resultats de cette experience.}


\section{Experience E06 --- Small Time-Frequency Patching}
\subsection{Choix et justification}
\textbf{Phase}: screening\\
\textbf{Seeds prevues}: 42

Tester un patching plus local en frequence pour enrichir la representation spectrale.

\paragraph{Hypothese attendue}
Qualite possiblement meilleure sur certains signaux, avec cout calcule plus eleve.

\subsection{Resultats}
\textit{A completer apres execution des runs.}

\subsection{Discussion des resultats}
\textit{A completer apres analyse des resultats de cette experience.}


\section{Experience E07 --- Small MAE mask 0.40}
\subsection{Choix et justification}
\textbf{Phase}: screening\\
\textbf{Seeds prevues}: 42

Ablation du ratio de masquage vers un regime plus conservateur.

\paragraph{Hypothese attendue}
Reconstruction plus facile, transfert potentiellement moins robuste.

\subsection{Resultats}
\textit{A completer apres execution des runs.}

\subsection{Discussion des resultats}
\textit{A completer apres analyse des resultats de cette experience.}


\section{Experience E08 --- Small MAE mask 0.75}
\subsection{Choix et justification}
\textbf{Phase}: screening\\
\textbf{Seeds prevues}: 42

Ablation du ratio de masquage vers un regime plus agressif.

\paragraph{Hypothese attendue}
Apprentissage plus difficile; possible gain de robustesse ou degradation selon capacite.

\subsection{Resultats}
\textit{A completer apres execution des runs.}

\subsection{Discussion des resultats}
\textit{A completer apres analyse des resultats de cette experience.}


\section{Experience E09 --- Final Top-1}
\subsection{Choix et justification}
\textbf{Phase}: final\\
\textbf{Seeds prevues}: 42, 123, 456, 789, 1024

Consolidation statistique de la meilleure variante screening (5 seeds).

\paragraph{Hypothese attendue}
Resultat principal du rapport.

\subsection{Resultats}
\textit{A completer apres execution des runs.}

\subsection{Discussion des resultats}
\textit{A completer apres analyse des resultats de cette experience.}


\section{Experience E10 --- Final Top-2}
\subsection{Choix et justification}
\textbf{Phase}: final\\
\textbf{Seeds prevues}: 42, 123, 456, 789, 1024

Consolidation statistique de la deuxieme meilleure variante screening.

\paragraph{Hypothese attendue}
Comparaison robuste contre Top-1.

\subsection{Resultats}
\textit{A completer apres execution des runs.}

\subsection{Discussion des resultats}
\textit{A completer apres analyse des resultats de cette experience.}


\section{Experience E11 --- Final Top-3}
\subsection{Choix et justification}
\textbf{Phase}: final\\
\textbf{Seeds prevues}: 42, 123, 456, 789, 1024

Consolidation statistique de la troisieme meilleure variante screening.

\paragraph{Hypothese attendue}
Completer le benchmark multi-modeles final.

\subsection{Resultats}
\textit{A completer apres execution des runs.}

\subsection{Discussion des resultats}
\textit{A completer apres analyse des resultats de cette experience.}


\section{Consolidation Top-3 (E09--E11)}
\subsection{Tableau final principal}
\textit{A completer: tableau mean $\pm$ std pour Top-1/Top-2/Top-3.}

\subsection{Visualisations prevues}
\begin{itemize}
\item Courbe WER par experience (E01 $\rightarrow$ E11).
\item Frontiere de Pareto WER vs runtime inference.
\item Frontiere de Pareto WER vs memoire GPU.
\end{itemize}
\textit{Inserer ici les figures et l'interpretation scientifique associee.}


\section{Limites et menaces a la validite}
\textit{A completer: limites compute/stockage, risques de variance, biais de selection des variantes, validite externe.}

\section{Conclusion}
\textit{A completer: synthese finale, reponse aux hypotheses, recommandations pour extension VoxPopuli et/ou TTS.}

\end{document}
